


\chapter{Literature Review}\label{chap:levantamento_literatura}


This chapter deepens the study of the Entanglement Theories and Speculative Design through a systematic investigation of the literature. Its goal is to map the state of the art, understand how different theoretical currents have been mobilized in Computing, and, in particular, analyze how these theories have been engaging (or not) with speculative design approaches.

The investigation is conducted through two complementary strategies: a bibliometric analysis followed by a systematic literature mapping (SLM). Together, these strategies provide an overview of the field and show that, although relevant intersections exist between Entanglement Theories and future-oriented design, these convergences remain largely underexplored when considered together. This finding reinforces the relevance of this dissertation’s contribution and prepares the ground for the framework presented in the following chapter.

\section{Contextualization}

The consolidation of a robust theoretical foundation on the Entanglement Theories within the field of Computing, particularly their articulation with Speculative Design, requires a methodological approach that goes beyond traditional literature reviews. Given the conceptual dispersion, epistemological fragmentation, and lack of methodological consensus observed in preliminary studies, a dual literature-review procedure was adopted, combining Bibliometric Analysis and SLM, as illustrated in Figure~\ref{fig:levantamento_literatura}.

\begin{figure}[htb!]
\centering
\caption{Stages of the literature review process.}
\includegraphics[width=1.0\textwidth]{FIGURAS/LEVANTAMENTO_LITERATURA.png}
\caption*{\textbf{Source:} Author’s elaboration.}
\label{fig:levantamento_literatura}
\end{figure}



The choice of distinct scopes for each stage stems from a strategic difference in objectives. The bibliometric analysis was conducted within the domain of IS, a field in which the Entanglement Theories have been historically mobilized and where there is a sufficiently robust body of publications to identify collaboration patterns, conceptual trends, and epistemological gaps. The SLM, in turn, focused on the field of Computing, with the aim of examining how these theories have been appropriated within Speculative Design, as well as in pedagogical, infrastructural, software development, and other related practices. This delimitation provides, on the one hand, a panoramic view of the theoretical debate in IS and, on the other, a situated analysis of applications in Computing, both of which are essential for supporting the objectives of this dissertation.

The bibliometric analysis follows principles of scientometrics \cite{hassan2010engaging}, understood as the quantitative and qualitative study of scientific production and communication. This approach makes it possible to identify patterns, relationships, and trends in the literature on Entanglement Theories in IS. The investigation revealed how these theories have been mobilized over the past decades. The data show a trajectory of growth, although marked by significant challenges such as the low rate of international collaboration (only 4.63\%) and the scarcity of integrative reviews. The results include thematic visualizations, co-occurrence networks, citation evolution, authorship patterns, and transnational collaborations.

The SLM stage was conducted with a qualitative and exploratory focus, aiming to understand how these theories have been appropriated, combined, and operationalized by studies that employ speculative practices in the field of Computing. The analysis of the 40 selected studies, following a protocol adapted from SEGRESS \cite{Kitchenham2022-SEGRESS}, revealed both innovative conceptual appropriations and recurring epistemological gaps, such as the superficial or merely illustrative use of theoretical approaches.

The complementarity of the two stages makes it possible to outline a more comprehensive picture of the state of the art and serves a strategic purpose: to ground the development of the framework that integrates the Entanglement Theories with Speculative Design (the central proposal of this dissertation), to clarify the relevance of the Entanglement Theories for speculative practices, and to guide the design of the Speculative Design tools associated with the framework.

\section{Bibliometric Analysis}

The Bibliometric Analysis was inspired by the framework proposed by \citeonline{ozturk2024design}. A summarized version of the results is presented here in order to avoid overwhelming the reader with excessive detail. Complete information, including methodological descriptions, visualizations, and extended analyses, can be found in Appendix \ref{apen:bibliometria}. Part of this material has also been consolidated into an article submitted to a high-impact international journal.

In this condensed version, the emphasis is placed on the main findings, particularly regarding the historical trajectory, the impact, and the collaboration networks surrounding these theories. The analysis examined aspects such as the volume and type of publications over time, academic recognition through citations, the role of central authors and their connections, and the convergence among theoretical references, influential authors, and recurring keywords. Additionally, a semantic analysis was conducted to identify thematic clusters and theoretical approaches that have gained prominence in recent debates.

\subsection{Research Method}

Studies published in English were collected, and the search strategy included both the central terms of these theories and their extensions, incorporating expressions such as ``Object-Oriented Ontology'' and ``Theory of Technological Mediation.'' Within the IS domain, terms representative of its thematic diversity were included, ranging from IT infrastructure and governance to sociotechnical design, HCI, data science, and artificial intelligence. This choice made it possible to capture approaches that span system development as well as the ethical and organizational implications of technology.

The bibliometric mapping was conducted in four distinct phases, beginning in September 2023 and updated periodically to follow the evolution of the literature. The final update occurred on 10 February 2025, when the complete dataset for this analysis was consolidated. Searches were performed in the Scopus and Web of Science databases. Document selection followed a structured three-stage process adapted from the model proposed by \citeonline{page2021prisma}.

In \textbf{stage 1}, a preliminary selection was conducted in the Scopus and Web of Science databases using search strings designed to capture the diversity of scientific publications related to the Entanglement Theories in IS and related fields. In \textbf{stage 2}, thematic filters specific to both databases were applied in order to restrict the results to the scope of IS and adjacent areas. The selected documents were then organized and processed through an evaluation of titles, abstracts, languages, and publication venues. An R script (see Appendix \ref{apen:script}) was developed to consolidate the data retrieved from the selected databases. Finally, in \textbf{stage 3}, additional documents were included and incomplete records were removed from the consolidated dataset, resulting in a total of 669 documents, which were exported to the Biblioshiny tool, in which descriptive and exploratory analyses of the bibliometric data were conducted.




\subsection{Data Analysis}

The data analysis of the Bibliometric Study sought to understand the evolution, impact, and thematic configuration of Entanglement Theories in IS. The first stage identified temporal publication trends, document types, and patterns of authorship. The subsequent stage examined academic impact by considering the average number of citations and the most influential articles. The authorship analysis mapped the main researchers and their collaboration networks using metrics such as citations normalized per year. The investigation also connected the main theoretical references to authors and keywords through three-dimensional visualizations. Finally, thematic evolution was explored through word clouds, co-occurrence networks, and temporal analyses.

\subsection{Results and Discussion of the Bibliometric Analysis}\label{sec:resultados_discussao_bibliometria}

The bibliometric analysis of the use of Entanglement Theories in IS reveals three central findings. First, it highlights the historical centrality of ANT as the predominant reference. Second, it identifies a set of limitations such as the low incidence of systematic reviews, which justifies conducting the SLM in the next stage of this study, and the limited international collaboration, indicating challenges for the theoretical consolidation and global articulation of knowledge in the field. Third, it identifies signs of conceptual reconfiguration with the recent rise of Agential Realism, Speculative Realism, and Postphenomenology, theories that expand the interpretive horizons in IS and point to a more critical and relational integration between humans, technologies, and sociotechnical practices. Each of these findings is discussed in detail below.

\subsubsection{Overview of Research on Entanglement Theories in IS}\label{subsec:panorama_geral}

To understand the development and consolidation of Entanglement Theories in IS, the evolution of the scientific output on the topic was examined over time (Table \ref{tab:resumo_analise_bibliometrica}). The analysis focused on publication volume, document types, and patterns of collaboration among authors.


\begin{table*}[ht!]
\centering
\scriptsize
\caption{General Overview of the Bibliometric Data}
\begin{tabular}{p{3.5cm}p{10.5cm}}

\toprule
\multicolumn{1}{l}{\textbf{Description}} & 
\multicolumn{1}{l}{\textbf{Quantity}}
\\

\cmidrule[0.4pt]{1-2}

\rowcolor{gray!10}
\multicolumn{2}{l}{\textbf{Main Information on the Data}}
\\

Period & January 1, 1992 to January 29, 2025
\\

Sources & 450 journals, books, and related outlets
\\

Documents & A total of 669 documents were analyzed\\

Annual Growth & Annual growth rate of 5\%
\\

Average Age & Average age of documents: 10.8 years
\\

Cit./Doc. & Mean of 17.23 citations per document
\\

References & A total of 4,943 cited references
\\

\cmidrule[0.4pt]{1-2}
\rowcolor{gray!10}
\multicolumn{2}{l}{\textbf{Authors}}
\\

Number of Authors & A total of 1,142 authors
\\

Single-Authored Docs. & A total of 210 documents written by a single author
\\

Co-Authors per Doc. & Mean of 2.26 co-authors per document
\\

International Co-Auth. \% & 4.634\% of documents involved international collaboration
\\

\cmidrule[0.4pt]{1-2}
\rowcolor{gray!10}
\multicolumn{2}{l}{\textbf{Document Types}}
\\

Articles & 304 articles published in scientific journals
\\

Conference Papers & 235 studies presented at conferences or workshops
\\

Reviews & 2 review papers
\\

Others & 128 editorial documents, notes, and special issue introductions
\\

\bottomrule
\end{tabular}
\label{tab:resumo_analise_bibliometrica}
\caption*{\textbf{Source:} Author’s elaboration}
\end{table*}


The bibliometric analysis, based on 669 documents published between 1992 and January 2025, reveals an average annual growth rate of 5\%, which indicates a continuous expansion of academic interest in the topic. Regarding scholarly collaboration, the data show that only 210 documents, approximately 31\%, are single-authored, while the overall mean is 2.26 co-authors per document. International collaboration remains limited, since only 4.63\% of the documents result from partnerships among authors from different countries. Concerning academic impact, the field exhibits an overall average of 17.23 citations per document. Although this is a reasonable indicator, it may be inflated by a small number of highly cited works. A total of 4,943 distinct references were cited, which suggests an extensive theoretical foundation.

A revealing aspect concerns the distribution of document types. Of the total analyzed, 304 documents, approximately 45\%, are journal articles, followed by 235 conference papers, corresponding to 35\%. Only two documents were classified as systematic reviews, which indicates a critical gap in the organization and assessment of the state of the art. This scarcity of review studies may hinder the consolidation of theoretical, conceptual, and methodological consensus. Additionally, the significant number of documents categorized as ``Others,'' consisting of 128 records that include editorial notes, introductions, and peripheral texts, may indicate a field that is still seeking legitimation and definition within the formal channels of science.

Overall, the current landscape portrays a promising and expanding field with increasing interdisciplinary interest. Nonetheless, it still faces important challenges, which include consolidating a more integrated theoretical identity, expanding international collaboration networks, and stabilizing its scientific output in a more continuous manner.

Despite the integrative potential of Entanglement Theories, combining them requires careful attention to the ontological tensions among approaches that adopt distinct views on entities, relations, and the constitution of reality. Such integration demands methodological caution and clarity regarding the ontological and epistemological premises of each theory in order to avoid collapsing fundamental differences into a generic notion of ``entanglement.''

\subsubsection{Key References and Their Articulation with Authors and Themes}\label{subsec:autres_temas}

This section analyzes how the most prominent external references cited in works related to Entanglement Theories in IS are connected to the most influential authors and to the recurrent themes identified through publication keywords. The aim is to understand how specific theoretical frameworks guide scientific production in the field and how these references are mobilized by different researchers and thematic domains.

Table \ref{tab:principais_referencias_citadas_local} provides an overview of the most frequently cited references in the analyzed corpus, highlighting the works that constitute the theoretical foundation for research on Entanglement Theories in IS.

 
\begin{table*}[ht!]
\centering
\scriptsize
\caption{Most Frequently Cited External References in the Analyzed Documents}
\rowcolors{2}{gray!10}{white}
\begin{tabular}{p{9.3cm}p{3.2cm}r}
\toprule
\textbf{Title} & \textbf{Author} & \textbf{Citations} \\
\cmidrule[0.4pt]{1-3}

Science in Action: How to Follow Scientists and Engineers Through Society & Bruno Latour & 27 \\

Sociological Review Monograph, P196 & Michel Callon and John Law & 21 \\

Power, Action and Belief: A New Sociology of Knowledge? & Michel Callon and John Law & 20 \\

A Sociology of Monsters: Essays on Power, Technology, and Domination & John Law & 20 \\

Actor-Network Theory and IS Research: Current Status and Future Prospects & Geoff Walsham & 17 \\

Notes on the Theory of the Actor-Network: Ordering, Strategy, and Heterogeneity & John Law & 16 \\

Why Has Critique Run Out of Steam? From Matters of Fact to Matters of Concern & Bruno Latour & 15 \\

Making Things Public: Atmospheres of Democracy & Isabelle Stengers, Bruno Latour, and Peter Weibel & 13 \\

Reassembling the Social: An Introduction to Actor-Network-Theory & Bruno Latour & 13 \\

We Have Never Been Modern & Bruno Latour & 12 \\

\bottomrule
\end{tabular}
\label{tab:principais_referencias_citadas_local}
\caption*{\textbf{Source:} Author’s elaboration.}
\end{table*}


ANT stands out historically as the main theoretical and philosophical foundation in investigations related to Entanglement Theories in IS. Among the ten most frequently cited titles, eight are authored or co-authored by Bruno Latour, Michel Callon, or John Law, who are the three central founders of ANT.

Although expressive, this result is also expected, given that ANT is the most established of the approaches analyzed. It reveals a clear preference for this theory within the scientific community and a form of paradigmatic stabilization in which ANT operates as an entry point for conceptualizing sociotechnical relations and the distributed agency of humans and nonhumans.

However, understanding the centrality of an approach requires more than recognizing its citation frequency. It also demands examining how it is mobilized in practice by researchers in the field. This involves going beyond citation counts to investigate how these references interact with the work of the most influential authors. In this context, Figure~\ref{fig:three_field_plot} maps the connections among leading authors, the references they most frequently cite, and the keywords associated with their publications. This visualization reveals how certain authors appropriate specific theoretical frameworks and articulate them with emerging themes, highlighting patterns of theoretical alignment and intellectual trajectories.

\begin{figure}[htb!]
\centering
\caption{Most Frequently Cited References by Leading Authors and the Most Prevalent Keywords}
\includegraphics[width=1.0\textwidth]{FIGURAS/TREE_FIELD.png}
\caption*{\textbf{Source:} Author’s elaboration using the Biblioshiny tool.}
\label{fig:three_field_plot}
\end{figure}



The diagram reveals a strong convergence among authors, references, and theoretical domains. Authors such as Atkinson C., Cecez-Kecmanovic D., Kautz K., and Elbanna A. maintain robust links to the works of Latour, Callon, and Law, indicating a solid theoretical affiliation with ANT. This adherence is also reflected in the keywords associated with their publications, including ``Actor-Network Theory,'' ``Information Systems,'' and ``Social Sciences Computing.''

While Cecez-Kecmanovic D. and Kautz K. associate the works of Latour and Callon predominantly with themes such as ``Information Systems'' and ``Social Sciences Computing,'' authors like Wickramasinghe N. and Elbanna A. broaden the thematic scope. They mobilize these same references to address topics such as information technologies, information use, philosophical perspectives, and even artificial intelligence. This expansion suggests that although a shared theoretical foundation exists, there are significant variations in how these theories are appropriated and articulated in relation to contemporary issues.

A particularly relevant case is the presence of \citeonline{walsham1997actor}, which appears among the most frequently cited references with 17 mentions. Walsham was one of the pioneers in articulating ANT with qualitative research in IS, emphasizing the analysis of interactions among technology, organizational context, and agency. His influence is evident among authors such as Iyamu T. and Brooks L., who appear to adopt this perspective to support investigations of a more interpretivist and sociotechnical nature.

Finally, the keyword ``Artificial Intelligence'' appears only modestly in the analyzed corpus, yet it represents a conceptually promising presence. Although it is not among the main terms associated with ANT in the work of the more traditional authors, it emerges in connection with researchers who also draw on the references of Latour and Callon. This suggests an effort to bring discussions on distributed agency and materiality into dialogue with the ethical and epistemic challenges posed by artificial intelligence. This emerging connection indicates that Entanglement Theories are beginning to be mobilized to problematize AI beyond its treatment as a mere technical tool, opening space for broader reflections on agency, ethics, and materiality in contemporary computational contexts.

\subsubsection{Thematic Evolution and Emerging Theoretical Approaches}\label{subsec:analise_tematica}

The emergence of AI as a recent theme reinforces the need to identify other conceptual trends in ascension, shifts in research focus, and the appearance of new theoretical perspectives over time. To address this, the evolution of themes was examined through an analysis of the keywords used in the publications. This approach makes it possible to understand how the field has been restructured in response to theoretical advances, interdisciplinary influences, and new challenges, revealing both continuities and innovations within the research domain. Figure \ref{fig:coocorrencia_palavras_chave} presents a keyword co-occurrence analysis based on the 669 documents examined.

\begin{figure}[htb!]
\centering
\caption{Keyword Co-occurrence Network}
\includegraphics[width=1.0\textwidth]{FIGURAS/CO_OCORRENCIAS.png}
\caption*{\textbf{Source:} Author’s elaboration using the Biblioshiny tool.}
\label{fig:coocorrencia_palavras_chave}
\end{figure}



The term ``Actor-Network Theory'' occupies the center of the network and functions as a connector among multiple domains. The green cluster shows the connection between ANT and more established applications in IS and organizational studies. Keywords such as ``information technology,'' ``project management,'' ``implementation,'' ``healthcare,'' ``developing countries,'' and ``case study'' reveal the instrumental and applied use of the theory in institutional, methodological, and practical contexts. This core represents a pragmatic line of appropriation directed toward the analysis of sociotechnical systems in real-world environments.

The blue cluster highlights a more conceptual axis oriented toward HCI, anchored in terms such as ``postphenomenology,'' ``phenomenology,'' ``research through design,'' and ``technological mediation.'' This network suggests connections between ANT and the phenomenological tradition, emphasizing interpretive methodologies and design-oriented critical approaches that are frequently associated with contemporary debates on ethics and technological practice.

The red cluster represents an ontologically and epistemologically dense strand structured around concepts associated with Agential Realism and Speculative Realism. Keywords such as ``agential realism,'' ``intra-action,'' ``entanglement,'' ``materiality,'' ``agency,'' ``ontology,'' and ``sociomateriality'' point to a deeper theoretical elaboration focused on the critique of the subject-object dualism and the reconfiguration of the notions of agency and reality. This group is also connected to contemporary topics such as ``ethics'' and ``artificial intelligence,'' indicating an effort to rethink the foundations of AI from perspectives that understand knowledge not as a representation of the world, but as something constructed within the interactions among people, technologies, and contexts. 

Although the co-occurrence analysis helps identify central themes and how they relate to one another, it does not show how these themes evolve over time. To address this limitation, a temporal trend analysis was conducted, as presented in Figure \ref{fig:trending_topics}.

\begin{figure}[htb!]
\centering
\caption{Trending Topics}
\includegraphics[width=1.0\textwidth]{FIGURAS/TRENDING_TOPICS.jpg}
\caption*{\textbf{Source:} Author’s elaboration using the Biblioshiny tool.}
\label{fig:trending_topics}
\end{figure}


In the early years of the historical series, terms such as ``interpretive research,'' ``structuration theory,'' ``actor,'' and ``implementation'' dominate the landscape, reflecting an initial interest in methodological foundations, organizational practices, and technology adoption processes.

Starting in 2010, a conceptual shift becomes visible, marked by the consolidation of ``actor-network theory'' as the structuring axis of the field. Its peak between 2012 and 2014 coincides with the rise of themes such as ``science and technology studies,'' ``case study,'' ``ontology,'' and ``materiality,'' indicating the expansion of ANT as a framework for understanding distributed agency, relations between humans and nonhumans, and the role of materiality in the constitution of IS.

Between 2015 and 2018, another theoretical shift takes place with the emergence of terms more closely aligned with contemporary Entanglement Theories and posthumanist approaches. Concepts such as ``sociomateriality,'' ``technology,'' ``agential realism,'' and ``materiality'' gain prominence and remain influential in subsequent years. This movement reveals a transition from the traditional organizational focus toward a deeper problematization of relational ontologies and the material conditions of action.

The recent emergence of ``entanglement'' as a central term constitutes the most significant indication of an ongoing attempt to consolidate Entanglement Theories in contemporary research. It is a concept that accurately captures the growing sensitivity to the inseparability of humans, artifacts, and sociotechnical practices, while also functioning as an articulating axis for multiple theoretical, methodological, and epistemological approaches oriented toward non-anthropocentric perspectives.

The nearly simultaneous appearance of ``artificial intelligence'' suggests a promising trend that brings debates on AI into closer dialogue with the perspectives offered by Entanglement Theories. The recent presence of ``human-computer interaction'' also draws attention, indicating a renewed convergence between these theories and critical studies in HCI.

This collective movement points to a field undergoing reconfiguration, in which new philosophical and ontological approaches begin to occupy a more prominent position, emerging as alternatives or complements to ANT. In this context, there is a visible effort to construct a more entangled perspective that integrates different theoretical approaches and enhances their sensitivity to the contemporary complexities of sociotechnical ecosystems.



\section{Systematic Literature Mapping}

The SLM sought to understand how Entanglement Theories have been mobilized by studies that employ Speculative Design. It aimed to identify trends and conceptual and methodological gaps, as well as to characterize the ways in which these theories are appropriated within the context of speculative practices.

\subsection{Research Method}

Six international databases were consulted: ACM Digital Library, IEEE Xplore, ISI Web of Science, ScienceDirect, Scopus, and SpringerLink. The search, conducted in January 2025, employed a combinatorial search string that intersected the axes of ``entanglement theories,'' ``speculative design,'' and ``technology.'' After applying inclusion and exclusion criteria, conducting full-text reading, and removing duplicates, 40 studies were selected. The process was carried out using the Parsif.al tool and validated by two reviewers, yielding a Cohen’s Kappa agreement index of 0.783.

\subsection{Data Analysis}

After the final selection of studies, the relevant data were extracted and analyzed. For each included article, information was collected regarding bibliographic data (year, authorship, country), the theories mobilized, the design practices described, the computational contexts of application, and the forms of conceptual appropriation.

The thematic categorization of the findings was guided both by these informational categories and by procedures inspired by Grounded Theory \cite{noble2016grounded}. Initially, one researcher conducted open coding of the studies, identifying recurring concepts (emerging themes) based on textual excerpts.

The resulting categories were refined through successive rounds of discussion among three researchers until a coherent set of patterns and conceptual gaps was consolidated.

\subsection{Results and Discussion of the Systematic Literature Mapping}

The analysis of the selected studies revealed different forms of articulation between Entanglement Theories and Speculative Design, highlighting both the potential for epistemological and critical advancement and the methodological and conceptual gaps that remain. The main findings are presented below and organized into three parts: (i) the characterization of the specific contributions of each theory to Speculative Design, (ii) the identification of gaps and emerging trends, and (iii) the analysis of the benefits and challenges associated with incorporating these theories into Technological Design practices.

\subsubsection{Synthesis of Related Work}

Table~\ref{tab:abordagens_e_DE} presents a synthesis of the ways in which different Entanglement Theories have been incorporated into Speculative Design within the context of Computing. The results show that each theory offers distinct contributions to the deepening and diversification of speculative practices in the field.


\begin{table*}[ht!]
\centering
\scriptsize
\caption{Exploration of Entanglement Theories through Speculative Design}
\arrayrulecolor{black}
\begin{tabular}{
    >{\raggedright\arraybackslash}p{3cm} 
    >{\raggedright\arraybackslash}p{5cm} 
    >{\raggedright\arraybackslash}p{5.5cm}
}

\toprule
\textbf{Relation} & \textbf{Exploration} & \textbf{Studies} \\
\cmidrule[0.4pt]{1-3}

\multirow{3}{*}{\shortstack{Postphenomenology and \\Speculative Design}}
  & Material speculation and co-speculation aligned with postphenomenological commitments to investigate human relations & \cite{forlano2023speculative,wakkary2022two} \\
  & Articulation of more expansive and pluralistic notions of relations & \cite{wakkary2022two} \\
  & Definition of agential and interactive affordances & \cite{kimura2023designing} \\
\arrayrulecolor{gray}
\cmidrule(l{0pt}r{0pt}){1-3}
\arrayrulecolor{black}

\multirow{3}{*}{\shortstack{OOO and \\Speculative Design}}
  & Design fiction used to explore object-oriented perspectives for designing systems such as the Internet of Things & \cite{akmal2020divination} \\
  & OOO-informed metaprobes as an extension of design probes & \cite{encinas2020metaprobes,kong2021neuromancer} \\
  & Questioning the ontological understanding of fictional objects & \cite{encinas2020metaprobes,coulton2024designing} \\
\arrayrulecolor{gray}
\cmidrule(l{0pt}r{0pt}){1-3}
\arrayrulecolor{black}

\multirow{3}{*}{\shortstack{ANT and\\ Speculative Design}}
  & Considering computational artifacts and algorithms as actors & \cite{bekker2023challenges,busboom2023futures,encinas2020metaprobes} \\
  & Research-through-Design (RtD) approaches combined with design futurization methods & \cite{busboom2023futures} \\
  & ANT used as metaphor-oriented design & \cite{encinas2020metaprobes,coulton2024designing} \\
\arrayrulecolor{gray}
\cmidrule(l{0pt}r{0pt}){1-3}
\arrayrulecolor{black}

\multirow{3}{*}{\shortstack{Agential Realism and\\ Speculative Design}}
  & Challenging dualist views and emphasizing the heightened agency of nonhuman actors in intra-actions with humans & \cite{bekker2023challenges,rozendaal2024enacting,kong2021neuromancer,lazar2021adopting} \\
  & Speculative enactments used to investigate experiences and social concerns related to alternative realities with technologies & \cite{kong2021neuromancer,rozendaal2024enacting,wakkary2022two,dorrenbacher2023intricacies} \\
  & Use of diffractive analysis & \cite{encinas2020metaprobes,forlano2023speculative,wakkary2022two} \\

\bottomrule
\end{tabular}
\label{tab:abordagens_e_DE}
\caption*{\textbf{Source:} Author’s elaboration.}
\end{table*}


%%%%%%%%%%%%%%%

Postphenomenology contributes by emphasizing technological mediations in the relations between humans and artifacts and by supporting speculative practices centered on lived experience and on the relational capacities of technologies. In the dialogue between ``Postphenomenology and Speculative Design,'' the TTM \cite{verbeek2005things} stands out as it supports practices of co-speculation and \textit{Worlding}\footnote{\textit{Worlding} is a concept that understands the world not as a fixed or stable entity but as a verb - a continuous process of world-making.} \cite{haraway2016staying}, highlighting the possibilities for action that emerge from co-constitutive relations between humans and technologies (agential affordances) \cite{kimura2023designing}.


The articulation between ``OOO and Speculative Design'' contributes by enabling the exploration of hidden layers of technical systems through design fiction and speculative tools. Added to this is the use of metaprobes\footnote{\textit{Metaprobes} are objects for metaphysical inquiry, or artifacts developed to reveal tensions, values, and imaginaries surrounding the speculated technologies.} \cite{encinas2020metaprobes}, speculative probes informed by OOO that operate as devices for ontological investigation by exploring the existence and agency of objects beyond their instrumental functions \cite{encinas2020metaprobes,coulton2024designing}. These practices aim to perform alternative ontologies and to instantiate new ways of existing in and with the world.


The presence of ANT in Speculative Design studies manifests at three levels: (i) the acknowledgment of the agency of computational artifacts \cite{bekker2023challenges}, which legitimizes the agency of nonhumans by considering artifacts and algorithms as actors within sociotechnical networks; (ii) its mobilization as a complementary theoretical structure for understanding sociotechnical relations and guiding speculative methods \cite{busboom2023futures}; and (iii) its metaphorical use as a narrative and critical framework \cite{encinas2020metaprobes, coulton2024designing}. Although conceptually relevant, ANT has been only minimally employed in its descriptive methodological dimension. For example, there is a notable absence of instruments native to ANT, such as controversy mapping\footnote{Controversy mapping is an ANT method that charts disputes and negotiations around still-unstable issues, revealing how different actants construct meanings, alliances, and conflicts over time.} \cite{venturini2021controversy}, which could deepen the analysis of sociotechnical networks.

Agential Realism, in turn, has been mobilized in speculative enactments and diffractive analysis\footnote{From a sociomaterial perspective, diffractive analyses promote the understanding that ideas, practices, and facts are the effects of entanglements and webs of relations among human and nonhuman actors.} \cite{barad2007meeting, forlano2023speculative} as ways to investigate how intra-actions between humans and technologies perform alternative realities, amplifying distinct perspectives on sociotechnical phenomena. Although Agential Realism has emerged as a reference, its theoretical complexity means that concepts such as intra-action and diffractive analysis are still being mobilized in partial or implicit ways. Nonetheless, a movement toward conceptual maturation can be observed, pointing to the gradual incorporation of this approach in a more robust manner.


\subsubsection{Gaps and Trends in the Articulation between Entanglement Theories and Speculative Design}

Table~\ref{tab:lacunas_tendencias_emaranhado_de} presents the main findings regarding the gaps and emerging trends in the articulation between Entanglement Theories and Speculative Design in the field of Computing.


\begin{table*}[ht!]
\centering
\scriptsize
\caption{Gaps and Trends in the Integration between Entanglement Theories and Speculative Design}
\arrayrulecolor{black}
\begin{tabular}{
    >{\raggedright\arraybackslash}p{1.0cm} 
    >{\raggedright\arraybackslash}p{7.2cm} 
    >{\raggedright\arraybackslash}p{5.3cm}
}

\toprule
\textbf{Type} & \textbf{Description} & \textbf{Studies} \\
\cmidrule[0.4pt]{1-3}

\multirow{6}{*}{Gaps}
  & Dependence on traditional Design/HCI conceptual repertoires & \cite{wakkary2022two} \\
  & Comprehensive Theoretical and Practical Engagement with AI & \cite{benjamin2023entoptic} \\
  & Integration of Ethics into Design Education & \cite{forlano2023speculative}, \cite{homewood2021tracing} \\
  & Design Methods for ``More-than-Human Bodies'' & \cite{homewood2021tracing} \\
  & Curricula Focused on Design Futures & \cite{forlano2023speculative} \\
\arrayrulecolor{gray}
\cmidrule(l{0pt}r{0pt}){1-3}
\arrayrulecolor{black}

\multirow{7}{*}{Trends}
  & Growing Integration of Entanglement Theories & \cite{dorrenbacher2023intricacies}, \cite{bekker2023challenges}, \cite{calderwood2024designing}, \cite{encinas2020metaprobes} \\
  & Design as Theoretical Inquiry & \cite{forlano2023speculative} \\
  & Focus on Ethics and Politics in Design & \cite{homewood2021tracing}, \cite{henry2024theorizing} \\
  & Use of Design Fiction and Speculative RtD & \cite{benjamin2023entoptic}, \cite{dorrenbacher2023intricacies}, \cite{encinas2020metaprobes} \\
  & Re-examination of Existing Models & \cite{henry2024theorizing} \\
  & Exploration of Co-creation with Nonhuman AI & \cite{tholander2023design} \\
  & Emphasis on Critical Educational Design & \cite{forlano2023speculative}, \cite{bekker2023challenges} \\

\bottomrule
\end{tabular}
\label{tab:lacunas_tendencias_emaranhado_de}
\caption*{\textbf{Source:} Author’s elaboration.}
\end{table*}

\citeonline{wakkary2022two} identify epistemological and terminological difficulties in incorporating philosophical foundations that lie outside the traditional repertoire of Design and HCI, which remains shaped by functionalist logics and usability-centered approaches. The study by \citeonline{benjamin2023entoptic} highlights a significant gap in the critical and comprehensive engagement with artificial intelligence technologies, particularly within design contexts. Although AI-based systems such as image-synthesis models are widely disseminated in everyday applications, their role in mediating human experiences of reality remains underexplored. Moreover, the authors observe that design research still treats AI predominantly as a tool or explanatory resource, neglecting its potential as speculative material and as an epistemologically active agent.

The study by \citeonline{homewood2021tracing} calls attention to gaps in design education with regard to the ethical formation of future professionals. Despite the increasing complexity of sociotechnical systems, educational practices continue to prioritize technical and anthropocentric competencies while overlooking approaches that foster critical thinking, ecological sensitivity, and social responsibility. \citeonline{forlano2023speculative} expand this critique by arguing that the most widespread branch of Speculative Design tends to generate hypothetical artifacts detached from concrete sociopolitical realities, neglecting the role of design as a situated and ethically responsible practice. The absence of literacies for ethical futures underscores the urgency of reconfiguring curricula in Computing, particularly in HCI, to include speculative, participatory, and post-human practices.


In contrast, several emerging trends seek to address these challenges. One of them is the growing integration of Entanglement Theories into HCI and Speculative Design through the adoption of posthumanist, multispecies, and performative approaches \cite{dorrenbacher2023intricacies, calderwood2024designing, bekker2023challenges}. This direction materializes in practices such as ``performative protocols,'' which embrace indeterminacy and recognize distributed agency as the ethical foundation for designing with AI \cite{calderwood2024designing}. The use of design fiction and speculative RtD approaches is also becoming consolidated as a means to problematize complex sociotechnical ecosystems, although authors such as \citeonline{henry2024theorizing} warn of their potential critical dilution when adopted in an instrumental manner. In response to this issue, there is a movement to reinterpret these models through post-anthropocentric perspectives.

In addition, there is an increasing push for a critical educational reorientation that values curricula more attentive to speculation, situated ethics, and the ontological complexity of sociotechnical systems \cite{forlano2023speculative, bekker2023challenges}. Finally, GenAI emerges as a transformative vector, acting both as a tool and as a co-authoring agent in the creative process, prompting a reconfiguration of the role of nonhumans in design \cite{tholander2023design}. 


\subsection{Conclusion of the Systematic Literature Mapping}

The combination of Entanglement Theories and Speculative Design in technological domains, especially in HCI, presents a distinct set of challenges, which are summarized in Table \ref{tab:limites_emaranhado_de}.


\begin{table*}[ht!]
\centering
\scriptsize
\caption{Challenges in Integrating Entanglement Theories and Speculative Design}
\rowcolors{2}{gray!10}{gray!2}
\begin{tabular}{p{6cm}p{9cm}}

\toprule
\multicolumn{1}{l}{\textbf{Challenge}} & 
\multicolumn{1}{l}{\textbf{Study}} \\
\cmidrule[0.4pt]{1-2}

\multicolumn{1}{l}{Complexity and Abstraction} & 
\cite{bekker2023challenges} \\

\multicolumn{1}{l}{Difficulty in Implementation and Evaluation} & 
\cite{kimura2022designing}, \cite{forlano2023speculative} \\

\multicolumn{1}{l}{Issues of Representation and Agency} & 
\cite{bekker2023challenges} \\

\multicolumn{1}{l}{Definition of Responsibility} & 
\cite{staahl2022making}, \cite{bekker2023challenges}, \cite{homewood2021tracing} \\

\multicolumn{1}{l}{Need for New Skills and Education} & 
\cite{forlano2023speculative}, \cite{homewood2021tracing} \\

\multicolumn{1}{l}{Risk of Marginalization} & 
\cite{forlano2023speculative}, \cite{henry2024theorizing} \\

\multicolumn{1}{l}{Managing Power Imbalances} & 
\cite{benjamin2023entoptic} \\

\bottomrule
\end{tabular}
\label{tab:limites_emaranhado_de}
\caption*{\textbf{Source:} Author’s elaboration.}
\end{table*}

One of the main challenges lies in the complexity and abstraction of these theories, which makes their operationalization difficult in educational or practical contexts \cite{bekker2023challenges}. Another challenge concerns how to represent the agency of more-than-human entities, such as algorithms, bacteria, forests, or sensors, without resorting to anthropocentric metaphors that project human qualities onto these agents or to practices that render them invisible in the design process \cite{bekker2023challenges}.

It is precisely at this point that the ethical challenge of ``defining responsibility'' emerges \cite{staahl2022making, bekker2023challenges, homewood2021tracing}. When imagining speculative futures, who is responsible for the visions being projected? Who is included or excluded from these futures? \citeonline{staahl2022making} draw attention to the risk that co-speculation may overlook the real impacts that certain worldviews can reinforce. With respect to the risk of marginalization, \citeonline{forlano2023speculative} emphasize that certain design practices, even when motivated by good intentions, may end up silencing voices and experiences that fall outside dominant normative patterns, such as those of people with disabilities, racialized populations, or historically excluded groups.

This problem becomes even more pronounced in the context of algorithmic technologies, which operate under claims of neutrality, yet may intensify inequalities and power asymmetries \cite{benjamin2023entoptic}. This requires design to reconsider its role as a mediator of multiple and asymmetric worlds, attending more carefully to nonhuman forms of agency.

In light of these challenges, \citeonline{forlano2023speculative, homewood2021tracing} argue for the reformulation of design curricula and pedagogical practices, with an emphasis on developing critical, speculative, and relational skills. However, this movement faces a significant obstacle: the difficulty of translating Entanglement Theories into educational practices without reducing their ontological complexity or instrumentalizing them in superficial ways \cite{bekker2023challenges}.

Addressing these challenges opens space for a series of benefits, which are presented in Table \ref{tab:beneficios_emaranhado_de}. 

\begin{table}[ht!]
\centering
\scriptsize
\caption{Benefits of Integrating Entanglement Theories and Speculative Design}
\rowcolors{2}{gray!10}{gray!2}
\begin{tabular}{p{6cm}p{8cm}}
\toprule
\textbf{Benefit} & \textbf{Justification} \\
\cmidrule[0.4pt]{1-2}

Theoretical Depth & 
Provides a robust ontological and epistemological foundation for understanding relations among humans, technologies, and the world as dynamic and non-binary intra-actions. \\

Expansion of the Design Scope & 
Expands the focus of design to include entire sociotechnical systems and nonhuman agents, opening new fields of speculation. \\

Enhanced Criticality and Reflection & 
Enables the exposure of values, assumptions, and often-invisible sociotechnical impacts, deepening the critical dimension of design. \\

Generation of Ethical and Political Insights & 
Helps formulate new ethical questions and political approaches by revealing distributed agency and interdependent relations. \\

Potential for Interdisciplinary Collaboration & 
Facilitates dialogue across areas such as philosophy, art, science, and technology, enriching design practices and perspectives. \\

Reframing Notions of Success and Failure & 
Reframes the understanding of failures as effects of complex sociomaterial tensions rather than as absences to be corrected. \\

\bottomrule
\end{tabular}
\label{tab:beneficios_emaranhado_de}
\caption*{\textbf{Source:} Author’s elaboration.}
\end{table}


First, there is a deepening of theoretical foundations, since design can rely on epistemologies that recognize the intra-action among humans, technologies, and environments. In this sense, the methodologies, tools, and design frameworks, as well as the intentions of facilitators and the dynamics enacted in speculative activities, actively participate in constituting meanings, shaping the projected sociotechnical ecosystems, and influencing the agencies involved. This recognition makes Speculative Design more attentive to the material-discursive conditions that shape processes of imagination, reflection, and creation.

Second, the scope of design expands to consider not only human users but also nonhuman agents, systems, and infrastructures \cite{forlano2023speculative, bekker2023challenges}. This broadened visibility of other agents refines the critical and reflective capacity of design by exposing implicit values, naturalized assumptions, and forms of silencing embedded in technological practices \cite{bekker2023challenges, forlano2023speculative, wakkary2022two}. It also reveals the normative layers that shape artifacts and computational systems by uncovering how certain worldviews are privileged while others are excluded. As a result, new ethical and political insights emerge \cite{homewood2021tracing, henry2024theorizing}, which are essential for imagining more just futures that are responsive to diverse experiences and forms of existence, particularly in contexts involving more-than-human agencies.

Moreover, this approach supports interdisciplinary collaboration by fostering productive dialogues among computing, philosophy, the arts, and the humanities \cite{calderwood2024designing, dorrenbacher2023intricacies}. It expands the heuristic potential of design and encourages the creation of hybrid methodologies capable of engaging with the sociotechnical complexity of contemporary ecosystems.

Finally, the articulation between Entanglement Theories and Speculative Design reframes the understanding of success and failure. Instead of seeking to correct errors according to traditional metrics, design begins to value the insights that emerge from sociomaterial tensions, ambiguity, and indeterminacy, which are constitutive elements of the realities in which it intervenes \cite{kimura2022designing, forlano2023speculative}.

\section{Final Considerations of the Chapter}

This chapter presented a literature review aimed at consolidating a robust theoretical and methodological foundation on Entanglement Theories and their articulation with Speculative Design in the field of Computing. The combination of the Bibliometric Analysis and the SLM proved to be strategic for understanding both the historical trajectory and the current state of academic discussions, bringing together broad quantitative data with a critical and in-depth qualitative interpretation.

The Bibliometric Analysis enabled the identification of key authors, collaboration networks, most cited works, and thematic trends that have shaped the development of Entanglement Theories in Information Systems over the past three decades. This overview revealed the centrality of ANT as a foundational reference, while also highlighting the recent emergence of approaches such as Agential Realism, Speculative Realism, and Postphenomenology, indicating a movement of theoretical reconfiguration in the field. Despite these advances, significant limitations were observed, including the low rate of international collaboration and the scarcity of integrative reviews. These gaps reinforce the need for more systematic and theoretically grounded investigations.

Complementing this, the SLM deepened the analysis by examining how these theories have been appropriated in studies that employ Speculative Design within computational contexts. The qualitative analysis of the 40 selected studies revealed both the potential for epistemological expansion and the methodological limitations in applying these approaches. Distinct forms of contribution from Entanglement Theories to speculative practices were identified, along with emerging trends such as the emphasis on ethical critique, co-speculation with nonhuman agents, and curricular reform aimed at addressing the complexities of contemporary sociotechnical systems.

The integration of bibliometric and systematic methods provides a more comprehensive view of the literature. Together, these approaches offer a mapping of the foundational milestones of Entanglement Theories, revealing how these foundations have been reinterpreted and applied in contemporary studies that investigate possible futures through Speculative Design. This body of related work therefore forms a continuous spectrum between foundational and recent contributions, supplying essential inputs for the formulation of the framework proposed in this dissertation.
